\exercisegroup

\begin{enumerate}
\def\labelenumi{\arabic{enumi})}
\tightlist
\item
  Let \(S\) be an arbitrary infinite set.
\end{enumerate}

\begin{enumerate}
\def\labelenumi{\alph{enumi})}
\tightlist
\item
  Show that the smallest cardinal of any total set in the normed space \(\mathcal{B}(\mathrm{S})\) of bounded mappings of S in \(\mathbf{R}\) (I, p.~4, Example) is equal to \(2^{\mathrm{Card}(\mathrm{S})}\) (consider the set of characteristic functions of subsets of S and note that there exists an enumerable set everywhere dense in \(\mathbf{R}\) ). b) Show that the smallest cardinal of any total set in the normed space \(\ell^1 (\mathrm{S})\) (I, p.~4, Example) is equal to \(\mathrm{Card}(\mathrm{S})\) .
\end{enumerate}

\begin{enumerate}
\def\labelenumi{\arabic{enumi})}
\setcounter{enumi}{1}
\item
  In the product topological vector space \(E = R^{N}\) over the field R, denote by \(e_{n} (n \in \mathbf{N})\) the elements of the canonical basis of the direct sum \(\mathbf{R}^{(\mathbf{N})}\) . Write \(a_{0} = e_{0}\) , \(a_{n} = e_{0} + (1/n)e_{n}\) for \(n \geqslant 1\) . Show that, for every integer \(n \geqslant 0\) , the \(a_{i}\) such that \(0 \leqslant i \leqslant n\) form a topologically independent family in E, but that the infinite family \((a_{n})_{n \geqslant 0}\) is not topologically independent. If M is the closed vector subspace \(Ra_{0}\) , the classes \(\dot{a}_{n}\) of the \(a_{n}\) in E/M form a topologically independent family (for \(n \geqslant 1\) ), but the closed vector subspace N generated by the \(a_{n}\) , with index \(n \geqslant 1\) , in E contains M.
\item
  Let E be a topological vector space over R, and f a homomorphism of the additive group of E in R. Show that if there exists a neighbourhood of 0 in E in which f is bounded, then f is a continuous linear form in E. This is so in particular when f is semi-continuous (lower or upper).
\item
  Denote by K the field R with the absolute value \(p(\xi) = |\xi|^{1/2}\) . Let E be the vector space over K of the real valued regulated functions defined over I = \([0, 1]\) , continuous on the right everywhere and zero at the point 1; show that on E the mapping \(x \mapsto \|x\| = \int_{0}^{1} |x(t)|^{1/2} dt\) is a norm. Show that for every function \(x \geqslant 0\) in E, there exists in E two functions \(x_{1} \geqslant 0\) , \(x_{2} \geqslant 0\) such that \(x = \frac{1}{2}(x_{1} + x_{2})\) and \(\|x_{1}\| = \|x_{2}\| = \frac{1}{\sqrt{2}} \|x\|\) . Deduce that every continuous linear form on E is identically zero.
\item
  Let K be a Hausdorff topological division ring of which the topology is locally retrobounded (GT, III, § 6, exerc. 22). Extend prop. 2 of I, p.~12 and th. 1 of I, p.~13 to topological vector spaces over K; similarly extend th. 2 of I, p.~13 and prop. 3 of I, p.~14 when K is also complete.
\item
  Let K be the topological division ring obtained by transferring the usual topology of \(Q^{2}\) to the field \(\mathbf{Q}(\sqrt{2})\) by the mapping \((x, y) \mapsto x + y\sqrt{2}\) .
\end{enumerate}

\begin{enumerate}
\def\labelenumi{\alph{enumi})}
\item
  Let \(\mathbf{E}\) be the set \(\mathbf{Q}(\sqrt{2})\) with its vector space structure over \(\mathbf{K}\) and with the topology induced by that of \(\mathbf{R}\) . Show that \(\mathbf{E}\) is a Hausdorff topological vector space, of dimension 1 on \(\mathbf{K}\) , but that it is not isomorphic to \(\mathbf{K}_s\) .
\item
  Let \(\mathbf{F}\) be the topological vector space \(\mathbf{E} \times \mathbf{E}\) over \(\mathbf{K}\) ; in \(\mathbf{F}\) , the hyperplane \(\mathbf{E} \times \{0\}\) is closed but there is no continuous linear form \(f\) on \(\mathbf{E} \times \mathbf{E}\) such that this hyperplane is given by the equation \(f(x) = 0\) .
\end{enumerate}

\begin{enumerate}
\def\labelenumi{\arabic{enumi})}
\setcounter{enumi}{6}
\item
  Let K be a valued division ring which is non-discrete and non-complete, let E be the topological vector subspace \(K + Ka\) of \(\hat{K}\) where \(a \notin K\) , and let F be the product space \(K \times E\) . In F, the subspace \(M = K \times \{0\}\) is closed and of codimension 2. Let N be the complementary subspace to M in F generated by the vectors (0, 1) and (1, a); show that F is not the direct topological sum of M and N.
\item
  Let \(p\) be a prime number, \(\mathbf{Q}_p\) the field of \(p\) -adic numbers (GT, III, § 6, exerc. 23). Let \(\mathrm{E}_0\) be the topological product space \(\mathbf{Q}_p \times \mathbf{R}\) ; if \(\mathbf{K}\) denotes the field \(\mathbf{Q}\) with the discrete topology, then \(\mathrm{E}_0\) is a topological vector space over \(\mathbf{K}\) . Let \(\mathbf{M}\) be the vector subspace formed by the elements \((r, r)\) where \(r\) varies in \(\mathbf{Q}\) ; further let \(\theta\) be an irrational number and \(\mathbf{N}\) the vector subspace formed by the elements \((0, r\theta)\) where \(r\) varies in \(\mathbf{Q}\) . Let \(\mathbf{E}\) be the subspace \(\mathbf{M} + \mathbf{N}\) of \(\mathrm{E}_0\) ; show that \(\mathbf{N}\) is a closed hyperplane in \(\mathbf{E}\) , but that there does not exist a complementary topological subspace to \(\mathbf{N}\) (note that \(\mathbf{M}\) is everywhere dense in \(\mathrm{E}_0\) ).
\item
  Let \(\mathbf{X}\) be a Hausdorff topological space, and let \(\mathbf{V}\) be a vector subspace of finite dimension \(n\) of the space \(\mathcal{C}(\mathbf{X};\mathbf{R})\) .
\end{enumerate}

\begin{enumerate}
\def\labelenumi{\alph{enumi})}
\item
  Show that there exist \(n\) pair-wise disjoint, open sets \(\mathrm{U}_i\) ( \(1 \leqslant i \leqslant n\) ) in \(\mathbf{X}\) , such that any function \(f \in \mathbf{V}\) which is identically zero in each of the \(\mathrm{U}_i\) , is identically zero in \(\mathbf{X}\) (use A, II, § 7.5, cor. 3).
\item
  Let \(x_{i} \in \mathbf{U}_{i}\) for \(1 \leqslant i \leqslant n\) . Deduce from a) that there exists a constant \(c > 0\) such that, for every function \(f \in \mathbf{V}\) , we have
\end{enumerate}

\[
\sup _ {x \in \mathrm{X}} | f (x) | \leqslant c \sum_ {i = 1} ^ {n} | f (x _ {i}) |.
\]

\begin{enumerate}
\def\labelenumi{\arabic{enumi})}
\setcounter{enumi}{9}
\tightlist
\item
  Let \(\mathbf{K}\) be a locally compact non-discrete valued division ring, and \(\mathbf{E}\) a left vector space of finite dimension over \(\mathbf{K}\) . Denote by \(\mathfrak{N}(\mathbf{E})\) the set of norms on \(\mathbf{E}\) , which is a subspace of the space \(\mathcal{C}(\mathbf{E};\mathbf{R})\) of mappings of \(\mathbf{E}\) , continuous (in the canonical topology), in \(\mathbf{R}\) .
\end{enumerate}

\begin{enumerate}
\def\labelenumi{\alph{enumi})}
\item
  When we give to \(\mathcal{C}(\mathrm{E};\mathbf{R})\) the compact convergence topology \(*\) (for which it is a Fréchet space), the set \(\mathfrak{N}(\mathrm{E})\) is closed in \(\mathcal{C}(\mathrm{E};\mathbf{R})\) , and locally compact.
\item
  Let \(p_0\) be an element of \(\mathfrak{N}(\mathrm{E})\) ; show that there exists a continuous mapping \((\lambda, p) \mapsto \pi_{\lambda}(p)\) of \([0, 1] \times \mathfrak{N}(\mathrm{E})\) in \(\mathfrak{N}(\mathrm{E})\) such that \(\pi_0(p) = p\) and \(\pi_1(p) = p_0\) for every \(p \in \mathfrak{N}(\mathrm{E})\) .
\end{enumerate}

\begin{enumerate}
\def\labelenumi{\arabic{enumi})}
\setcounter{enumi}{10}
\tightlist
\item
  With the hypotheses of I, p.~23, exerc. 7 show that if K and E are complete then every closed subspace of E has a topological complement (proceed as in a), loc. cit.).
\end{enumerate}

¶ 12) Let K be a locally compact valued division ring whose absolute value is non-discrete and ultrametric. We call a norm on the left vector space E over K an ultranorm if it satisfies the ultrametric inequality (II, p.~2).

\begin{enumerate}
\def\labelenumi{\alph{enumi})}
\tightlist
\item
  Let E be a finite dimensional left vector space over K, let \(\alpha\) be an ultranorm on E and H a hyperplane in E given by the equation \(\langle x, a^{*} \rangle = 0\) . Show that there exists a point \(x_{0} \in E\) at which the function \(x \mapsto |\langle x, a^{*} \rangle| / \alpha(x)\) attains its upper bound in \(E \setminus \{0\}\) ; show that then
\end{enumerate}

\[
\alpha (x) = \sup \left(\alpha \left(x - \frac {\langle x , a ^ {*} \rangle}{\langle x _ {0} , a ^ {*} \rangle} x _ {0}\right), \frac {| \langle x , a ^ {*} \rangle |}{| \langle x _ {0} , a ^ {*} \rangle |} \alpha (x _ {0})\right).
\]

Deduce that there exists a basis \((a_{i})\) of E and a family \((r_i)\) of real numbers \(>0\) such that, for all \(x = \sum_{i} \xi_{i} a_{i}\) we have \(\alpha(x) = \sup_{i} (r_{i} |\xi_{i}|)\) . We say that \(\alpha\) is in the standard form relative to the basis \((a_{i})\) .

\begin{enumerate}
\def\labelenumi{\alph{enumi})}
\setcounter{enumi}{1}
\tightlist
\item
  Let \(\alpha^{*}\) be the norm on E* the dual of E canonically associated with \(\alpha\) by
\end{enumerate}

\[
\alpha^ {*} (x ^ {*}) = \sup _ {x \neq 0} | \langle x, x ^ {*} \rangle | / \alpha (x);
\]

it is an ultranorm. Show that for all \(x_0 \neq 0\) in \(\mathbf{E}\) , there exists \(x_0^* \in \mathbf{E}^*\) such that \(\alpha(x_0) = |\langle x_0, x_0^* \rangle| / \alpha^*(x_0^*)\) .

\begin{enumerate}
\def\labelenumi{\alph{enumi})}
\setcounter{enumi}{2}
\tightlist
\item
  Let \(\alpha, \beta\) be any two ultranorms on E. Show that there exists a basis of E such that relative to this basis \(\alpha\) and \(\beta\) are both of the standard form (consider a point \(x_0 \in \mathrm{E} \setminus \{0\}\) at which \(\alpha / \beta\) attains its maximum; then use \(b\) ) and proceed by induction on \(\dim \mathrm{E}\) ).
\end{enumerate}

\(d)\) Let \(\mathfrak{N}_0(\mathrm{E})\) , the set of ultranorms on \(\mathrm{E}\) , be considered as a subspace of \(\mathfrak{N}(\mathrm{E})\) (exerc. 10). Show that \(\mathfrak{N}_0(\mathrm{E})\) is closed in \(\mathfrak{N}(\mathrm{E})\) . Let \(\alpha_0\) be an element of \(\mathfrak{N}_0(\mathrm{E})\) ; for each \(\alpha \in \mathfrak{N}_0(\mathrm{E})\) and for \(0 \leqslant t \leqslant 1\) , let \(\mathrm{P}_{\alpha}(t)\) be the set of \(\beta \in \mathfrak{N}_0(\mathrm{E})\) such that \(\beta(x) \leqslant \alpha_0(x)^{1 - t}\alpha(x)^t\) for all \(x \in \mathrm{E}\) . Show that \(\mathrm{P}_{\alpha}(t)\) is not empty and that \(\pi_t^\alpha = \sup \mathrm{P}_{\alpha}(t)\) is an ultranorm. Further, the mapping \((t, \alpha) \mapsto \pi_t^\alpha\) of \([0, 1] \times \mathfrak{N}_0(E)\) in \(\mathfrak{N}_0(E)\) is continuous and such that \(\pi_0^\alpha = \alpha_0\) and \(\pi_1^\alpha = \alpha\) (use \(c\) ).

* e) Let A be the ring of the absolute value of K, m its maximal ideal such that \(k = A/m\) is a finite field with q elements (CA, VI, § 5, No.~1, prop. 2). For every ultranorm \(\alpha\) on E, the image \(X_{\alpha}\) of the set of values of \(\log \alpha(x)\) for \(x \in E \setminus \{0\}\) under the canonical mapping in the quotient group \(\mathbf{R}/(\mathbf{Z}, \log q)\) is a finite set having at most \(n = \dim E\) elements (use a); the number of these elements is denoted by \(r(\alpha)\) and called the rank of \(\alpha\) . Show that r is a lower semi-continuous mapping of \(\mathfrak{N}_{0}(E)\) in N and that the set \(\mathfrak{N}_{0}'(E)\) of the \(\alpha\) for which \(r(\alpha) = n\) is open and everywhere dense in \(\mathfrak{N}_{0}(E)\) (use a) and c)).

\begin{enumerate}
\def\labelenumi{\alph{enumi})}
\setcounter{enumi}{5}
\tightlist
\item
  Suppose that \(r(\alpha) = n\) ; let \((a_i)\) be a basis of E relative to which \(\alpha\) is of the standard form; show that there exists a neighbourhood V of \(\alpha\) in \(\mathfrak{N}_0'(\mathrm{E})\) such that every \(\beta \in V\) has the standard form relative to \((a_i)(\text{use } b)\) ; deduce that there is a neighbourhood \(W \subset V\) of \(\alpha\) homeomorphic to an open set in \(R^n\) .
\end{enumerate}

\(g)\) For every basis \((a_i)\) of E, show that the set of ultranorms \(\alpha\) that have a standard form relative to \((a_i)\) is closed in \(\mathfrak{N}_0^{\prime}(\mathrm{E})\) . Deduce that if \(\alpha \in \mathfrak{N}_0^{\prime}(\mathrm{E})\) has a standard form relative to \((a_i)\) the same is true of all elements of the connected component containing \(\alpha\) in \(\mathfrak{N}_0^{\prime}(\mathrm{E})\) .

¶ 13) * We keep the general hypotheses and the notations of exerc. 12.

\begin{enumerate}
\def\labelenumi{\alph{enumi})}
\item
  Let \(\mathbf{L}\) be a free sub-A-module of \(\mathbf{E}\) of dimension \(n = \dim \mathbf{E}\) . For all \(x \in \mathbf{E} \setminus \{0\}\) , the set of the \(a \in \mathbf{A}\) such that \(ax \in \mathbf{L}\) is a fractional ideal of \(\mathbf{K}\) of the form \(\mathfrak{m}^h\) ( \(h\) a positive or negative integer); putting \(\alpha(x) = q^h\) and \(\alpha(0) = 0\) , show that \(\alpha\) is an ultranorm on \(\mathbf{E}\) . It is said to be associated with the free A-module \(\mathbf{L}\) .
\item
  Conversely, if \(\alpha\) is an ultranorm on E, the set \(L_{\alpha}\) of the \(x \in E\) such that \(\alpha(x) \leqslant 1\) is a free A-module of dimension \(n\) . If \([\alpha]\) is the norm associated with \(L_{\alpha}\) , we have \(\alpha \leqslant [\alpha] \leqslant q\alpha\) , and \([\alpha]\) is the lower bound of the norms associated with free A-modules and which are \(\geqslant \alpha\) . We have \([q\alpha] = q \cdot [\alpha]\) , and \(\alpha(x) = \inf(q^{-t}[q^t\alpha](x))\) for all \(x \in E\) , where \(t\) varies in the interval {[}0, 1{]}. Further, the function \(t \mapsto [q^t\alpha](x)\) is left continuous in this interval.
\item
  With the same notations, show that for \(0 \leqslant t \leqslant 1\) , there are at most n distinct ultranorms among the \([q^{t}\alpha]\) . Conversely, let L be the set of ultranorms associated with the free A-modules of dimension n, and let \((\alpha_{t})_{0 \leqslant t \leqslant 1}\) be an increasing family of ultranorms of L such that \(\alpha_{1} = q\alpha_{0}\) . Show that there exists a basis of E relative to which all the \(\alpha_{t}\) have the standard form (if \(u \in A\) is an element of valuation 1, and \(L_{t}\) the free A-module of the \(x \in E\) such that \(\alpha_{t}(x) \leqslant 1\) , consider the vector spaces \(L_{t}/uL_{0}\) on k). Deduce further, that if, for all \(x \in E\) , \(t \mapsto \alpha_{t}(x)\) is left-continuous in \([0, 1]\) then there exists a unique ultranorm \(\alpha\) such that \(\alpha_{t} = [q^{t}\alpha]\) for all \(t \in [0, 1]\) .
\end{enumerate}

\(d\) ) The linear group \(\mathbf{GL}(\mathrm{E})\) operates continuously in \(\mathfrak{N}_0(\mathrm{E})\) ; show that it operates properly. For all \(\alpha \in \mathfrak{N}_0(\mathrm{E})\) , the stabiliser \(S_{\alpha}\) of \(\alpha\) in \(\mathbf{GL}(\mathrm{E})\) is the intersection of the stabilisers of the \([q^t\alpha]\) for \(0\leqslant t\leqslant 1\) ; deduce that \(S_{\alpha}\) is an open compact subgroup of \(\mathbf{GL}(\mathrm{E})\) , and hence that the orbit of each \(\alpha \in \mathfrak{N}_0(\mathrm{E})\) is a closed, discrete subspace of \(\mathfrak{N}_0(\mathrm{E})\) .

\begin{enumerate}
\def\labelenumi{\alph{enumi})}
\setcounter{enumi}{4}
\tightlist
\item
  For every ultranorm \(\alpha \in \mathfrak{N}_0(\mathrm{E})\) consider the decreasing sequence of the dimensions of the vector \(k\) -spaces \(\mathrm{L}_t / u\mathrm{L}_0\) , where \(\mathrm{L}_t\) is the A-module of the \(x \in \mathrm{E}\) such that \([q^t\alpha](x) \leqslant 1\) , and \(t\) varies from 0 to 1; we call this sequence, the sequence of invariants of \(\alpha\) . In order that \(\alpha\) and \(\beta\) belong to the same orbit in \(\mathfrak{N}_0(\mathrm{E})\) , it is necessary and sufficient that \(\mathrm{X}_{\alpha} = \mathrm{X}_{\beta}\) (exerc. 12, e)) and that the sequence of the invariants of \(\alpha\) and of \(\beta\) should be the same (use exerc. 12, b)). f) Deduce from e) that the space of the orbits \(\mathfrak{N}_0(\mathrm{E}) / \mathbf{GL}(\mathrm{E})\) is isomorphic with the space of the orbits \(\mathbf{T}^n / \mathfrak{S}_n\) , where the symmetric group operates on the right on \(\mathbf{T}^n\) by \((z_1, ..., z_n) \mapsto (z_{\sigma(1)}, ..., z_{\sigma(n)})\) . *
\end{enumerate}

\begin{enumerate}
\def\labelenumi{\arabic{enumi})}
\setcounter{enumi}{13}
\item
  Generalize the results of No.~2 and No.~3 to topological vector spaces E over a discrete division ring K, such that there exists a fundamental system of balanced neighbourhoods of 0 in E (i.e.~of neighbourhoods V such that K.V = V).
\item
  Let \(\mathbf{E}\) be a normed space of finite dimension \(n\) over \(\mathbf{R}\) or \(\mathbf{C}\) . Ascribe to the dual \(\mathbf{E}^*\) , the norm defined by \(\| x^* \| = \sup_{\| x \| \leqslant 1} |\langle x, x^* \rangle|\) (GT, X, § 3.2). Show that there exists a basis \((e_i)\)
\end{enumerate}

of E such that, if \((e_{i}^{*})\) is the dual basis, we have \(\|e_{i}\| = \|e_{i}^{*}\| = 1\) for all i. (Let \((a_{i})\) be a basis of E formed of vectors of norm 1; consider, the determinant \(\det(\xi_{ij})\) for each system of n vectors \(x_{i} = \sum_{j} \xi_{ij} a_{j}\) of norm 1, and consider such a system for which the absolute value of this determinant is maximal.)

